\subsection{How to solve a general penalized MLE problem}

% slides P68

The optimization problem is
\[
\hat\beta = \argmin_\beta \ell(\beta) + \lambda p(\beta),
\]
where $\ell(\beta)$ is a sufficiently smooth loss function and $p(\beta)$ is a general penalty function.

We develop an iterative algorithm in steps as follows. The idea is that we first focus on an optimization procedure on $\ell(\beta)$, then consider the additional penalty term $\lambda p(\beta)$ in the optimization procedure.
\begin{enumerate}
    \item \textbf{Step 1:} Apply linear quadratic approximation (LQA) on $\ell(\beta)$ with a second-order Taylor expansion around a fixed $\tilde{\beta}$. Denote the LQA function by $\ell_Q(\beta)$. 
    \[
    \argmin_\beta \ell_Q(\beta) + \lambda p(\beta)
    \]
    \item \textbf{Step 2:} We first develop the algorithm to solve the unpenalized loss:
    \[
    \argmin_\beta \ell_Q(\beta).
    \]
    This can be solved by the Majorization-Minimization (MM) algorithm. See Page 69.
\end{enumerate}


% logistic regr example.
% For most choices of $p$ (even if not differentiable), you can rewrite $\ell_Q(\beta)$ in the form of (generalized) least squares. Then apply coordinate descent algorithm. as done in Lasso. The additional L1 penalty term is handled by the additional soft-thresholding step. This is a special case of the proximal gradient descent algorithm.